%% ----------------------------------------------------------------------------
%%  §10.4  La formula di Taylor con il resto di Lagrange
%% ----------------------------------------------------------------------------
\section{La formula di Taylor con il resto di Lagrange}
\label{sec:resto-lagrange}

Il resto integrale è esatto, ma per stimarlo non occorre calcolare l'integrale: basta
sapere tra quali valori varia $f^{(n+1)}$. Il peso $\frac{(x-t)^{n}}{n!}$, infatti, non cambia segno
nell'intervallo di integrazione, e per una media pesata con peso di segno costante vale
un analogo del teorema della media (\cref{teo:media}).

\begin{lemma}[Teorema della media pesata]\label{lem:media-pesata}
\index{teorema!della media pesata}
Siano $g$ e $w$ continue in $[\alpha,\beta]$, con $w(t)\ge0$ per ogni $t\in[\alpha,\beta]$. Esiste un punto
$\xi\in[\alpha,\beta]$ tale che
\begin{equation}\label{eq:media-pesata}
  \int_{\alpha}^{\beta}g(t)\,w(t)\,dt=g(\xi)\int_{\alpha}^{\beta}w(t)\,dt .
\end{equation}
\end{lemma}

\begin{schema}
\Ipotesi{$g$, $w$ continue in $[\alpha,\beta]$; \ $w\ge0$.}
\Tesi{$\esiste\xi\in[\alpha,\beta]$ : $\int_{\alpha}^{\beta}g\,w=g(\xi)\int_{\alpha}^{\beta}w$.}
\end{schema}

\begin{strategia}
Si stringe $g$ tra il suo minimo $m$ e il suo massimo $M$ e si moltiplica per $w\ge0$, che
conserva le disuguaglianze; integrando si trova che il rapporto
$\frac{\int gw}{\int w}$ è compreso tra $m$ e $M$, e il teorema dei valori intermedi fornisce il
punto $\xi$. Va trattato a parte il caso $\int w=0$, in cui il rapporto non ha senso.
\end{strategia}

\begin{dimostrazione}
Per il teorema di Weierstrass (\cref{teo:weierstrass}) esistono
$m=\min_{[\alpha,\beta]}g$ e $M=\max_{[\alpha,\beta]}g$. Poiché $w\ge0$, da $m\le g(t)\le M$ segue
$m\,w(t)\le g(t)\,w(t)\le M\,w(t)$ per ogni $t$, e per il confronto tra integrali
(\cref{teo:confronto-integrali}), essendo $gw$ continua e quindi integrabile,
\begin{equation}\label{eq:media-pesata-stima}
  m\int_{\alpha}^{\beta}w(t)\,dt\le\int_{\alpha}^{\beta}g(t)\,w(t)\,dt\le M\int_{\alpha}^{\beta}w(t)\,dt .
\end{equation}
Se $\int_{\alpha}^{\beta}w=0$, la \eqref{eq:media-pesata-stima} dà $\int_{\alpha}^{\beta}gw=0$, e la tesi vale con
qualunque $\xi$. Altrimenti $\int_{\alpha}^{\beta}w>0$, e dividendo si ottiene
$m\le\frac{\int_{\alpha}^{\beta}gw}{\int_{\alpha}^{\beta}w}\le M$. Per il secondo teorema dei valori intermedi
(\cref{teo:valori-intermedi-2}) la funzione continua $g$ assume in $[\alpha,\beta]$ tutti i valori
compresi tra $m$ e $M$: esiste dunque $\xi$ con $g(\xi)=\frac{\int gw}{\int w}$, che è la tesi.
\end{dimostrazione}

Con $w\equiv1$ si ritrova il teorema della media. Applichiamo ora il lemma al resto
integrale.

\begin{teorema}[Formula di Taylor con il resto di Lagrange]\label{teo:resto-lagrange}
\index{formula di Taylor!con il resto di Lagrange}
Sia $f$ derivabile $n+1$ volte in $[a,b]$, con derivata $f^{(n+1)}$ continua, e sia
$x_{0}\in[a,b]$. Per ogni $x\in[a,b]$ esiste un punto $x_{1}$ compreso tra $x_{0}$ e $x$ tale che
\begin{equation}\label{eq:resto-lagrange}
  R_{n}(x)=\frac{f^{(n+1)}(x_{1})}{(n+1)!}\,(x-x_{0})^{n+1} ,
\end{equation}
cioè
\[
  f(x)=\sum_{k=0}^{n}\frac{f^{(k)}(x_{0})}{k!}(x-x_{0})^{k}+\frac{f^{(n+1)}(x_{1})}{(n+1)!}(x-x_{0})^{n+1} .
\]
\end{teorema}

\begin{schema}
\Ipotesi{$f\in C^{n+1}([a,b])$; \ $x_{0},x\in[a,b]$.}
\Tesi{$\esiste x_{1}$ tra $x_{0}$ e $x$ : \ $R_{n}(x)=\frac{f^{(n+1)}(x_{1})}{(n+1)!}(x-x_{0})^{n+1}$.}
\end{schema}

\begin{strategia}
Il resto è l'integrale \eqref{eq:resto-integrale} di $f^{(n+1)}$ contro il peso
$\frac{(x-t)^{n}}{n!}$, che ha segno costante tra $x_{0}$ e $x$. Il \cref{lem:media-pesata} porta
fuori dall'integrale $f^{(n+1)}$, calcolata in un punto intermedio, e l'integrale del peso
si calcola esplicitamente. Se $x<x_{0}$ occorre prima riportare l'integrale a un intervallo
orientato e un peso non negativo, controllando i segni.
\end{strategia}

\begin{dimostrazione}
Se $x=x_{0}$ si ha $R_{n}(x_{0})=0$ e la tesi vale con $x_{1}=x_{0}$.

\emph{Caso $x>x_{0}$.} Il peso $w(t)=\frac{(x-t)^{n}}{n!}$ è continuo e non negativo in
$[x_{0},x]$, e
\begin{equation}\label{eq:integrale-peso}
  \int_{x_{0}}^{x}\frac{(x-t)^{n}}{n!}\,dt=\Bigl[-\frac{(x-t)^{n+1}}{(n+1)!}\Bigr]_{t=x_{0}}^{t=x}
  =\frac{(x-x_{0})^{n+1}}{(n+1)!} .
\end{equation}
Per il \cref{teo:resto-integrale} e il \cref{lem:media-pesata} con $g=f^{(n+1)}$, esiste
$x_{1}\in[x_{0},x]$ tale che
\[
  R_{n}(x)=\int_{x_{0}}^{x}\frac{(x-t)^{n}}{n!}\,f^{(n+1)}(t)\,dt
  =f^{(n+1)}(x_{1})\,\frac{(x-x_{0})^{n+1}}{(n+1)!} .
\]

\emph{Caso $x<x_{0}$.} Scambiando gli estremi e scrivendo $(x-t)^{n}=(-1)^{n}(t-x)^{n}$,
\[
\begin{aligned}
  R_{n}(x)&=-\int_{x}^{x_{0}}\frac{(x-t)^{n}}{n!}\,f^{(n+1)}(t)\,dt\\
  &=(-1)^{n+1}\int_{x}^{x_{0}}\frac{(t-x)^{n}}{n!}\,f^{(n+1)}(t)\,dt .
\end{aligned}
\]
Ora il peso $\frac{(t-x)^{n}}{n!}$ è non negativo in $[x,x_{0}]$, e il suo integrale vale
$\frac{(x_{0}-x)^{n+1}}{(n+1)!}$; per il \cref{lem:media-pesata} esiste $x_{1}\in[x,x_{0}]$ tale che
\[
  R_{n}(x)=(-1)^{n+1}f^{(n+1)}(x_{1})\,\frac{(x_{0}-x)^{n+1}}{(n+1)!}
  =f^{(n+1)}(x_{1})\,\frac{(x-x_{0})^{n+1}}{(n+1)!} ,
\]
perché $(-1)^{n+1}(x_{0}-x)^{n+1}=(x-x_{0})^{n+1}$.
\end{dimostrazione}

\begin{osservazione}\label{oss:lagrange-n0}
Per $n=0$ la \eqref{eq:resto-lagrange} si riduce a
$f(x)-f(x_{0})=R_{0}(x)=f'(x_{1})\,(x-x_{0})$: è il teorema di Lagrange
(\cref{teo:lagrange}), di cui la formula di Taylor con il resto di Lagrange è dunque
una generalizzazione. Si noti che il teorema di Lagrange richiede soltanto la
derivabilità di $f$, e non la continuità di $f'$: l'ipotesi di continuità di $f^{(n+1)}$
nel \cref{teo:resto-lagrange} è dovuta al metodo dimostrativo, che passa per il resto
integrale, e non è necessaria. Lo mostriamo ora, generalizzando il procedimento usato
nel \cref{lem:taylor-lagrange} del \cref{cap:piu-variabili}.
\end{osservazione}

\begin{teorema}[Resto di Lagrange, ipotesi minime]\label{teo:resto-lagrange-generale}
Sia $f$ derivabile $n+1$ volte in un intervallo aperto $I$, e siano $x_{0},x\in I$ con
$x\neq x_{0}$. Esiste un punto $x_{1}$ strettamente compreso tra $x_{0}$ e $x$ per cui vale la
\eqref{eq:resto-lagrange}.
\end{teorema}

\begin{schema}
\Ipotesi{$f$ derivabile $n+1$ volte in $I$ aperto (nessuna ipotesi di continuità su
$f^{(n+1)}$); \ $x_{0}\neq x$ in $I$.}
\Tesi{$\esiste x_{1}$ strettamente tra $x_{0}$ e $x$ : \ $R_{n}(x)=\frac{f^{(n+1)}(x_{1})}{(n+1)!}(x-x_{0})^{n+1}$.}
\end{schema}

\begin{strategia}
Confrontiamo il resto $F(t)=R_{n}(t)$ con la funzione $G(t)=(t-x_{0})^{n+1}$, che ha la
stessa «struttura»: entrambe si annullano in $x_{0}$ insieme alle loro prime $n$
derivate. Il rapporto $\frac{F(x)}{G(x)}$ è quindi un rapporto di incrementi, a cui si applica il
teorema di Cauchy; il risultato è un rapporto di derivate prime, che è ancora un
rapporto di incrementi tra $x_{0}$ e un punto intermedio, e così via. Dopo $n+1$ passi
compaiono le derivate di ordine $n+1$: $F^{(n+1)}=f^{(n+1)}$ e $G^{(n+1)}=(n+1)!$. È lo
stesso schema della dimostrazione del teorema di L'Hôpital, percorso al finito.
\end{strategia}

\begin{dimostrazione}
Supponiamo $x>x_{0}$ (il caso $x<x_{0}$ è identico, con gli intervalli $[\xi,x_{0}]$ al posto
di $[x_{0},\xi]$). Siano $F(t)=f(t)-T_{n}(t)$ e $G(t)=(t-x_{0})^{n+1}$ per $t\in I$. Poiché $T_{n}$ è
un polinomio di grado $\le n$, $F$ è derivabile $n+1$ volte in $I$, con
$F^{(n+1)}=f^{(n+1)}$; per il \cref{lem:derivate-taylor},
$F^{(k)}(x_{0})=f^{(k)}(x_{0})-T_{n}^{(k)}(x_{0})=0$ per $k=0,1,\dots,n$. Inoltre
\[
  G^{(k)}(t)=\frac{(n+1)!}{(n+1-k)!}\,(t-x_{0})^{n+1-k} \qquad (k=0,1,\dots,n+1),
\]
sicché $G^{(k)}(x_{0})=0$ per $k\le n$, $G^{(k)}(t)\neq0$ per $t\neq x_{0}$, e $G^{(n+1)}\equiv(n+1)!$.

Poniamo $\xi_{0}=x$ e costruiamo per ricorrenza punti $x_{0}<\xi_{n+1}<\xi_{n}<\dots<\xi_{1}<\xi_{0}$ tali
che
\begin{equation}\label{eq:catena-cauchy}
  \frac{F^{(k-1)}(\xi_{k-1})}{G^{(k-1)}(\xi_{k-1})}=\frac{F^{(k)}(\xi_{k})}{G^{(k)}(\xi_{k})} \qquad (k=1,\dots,n+1) .
\end{equation}
Dato $\xi_{k-1}\in(x_{0},x]$, con $k\le n+1$, le funzioni $F^{(k-1)}$ e $G^{(k-1)}$ sono derivabili in
$I$ (perché $k-1\le n$), dunque continue in $[x_{0},\xi_{k-1}]$ e derivabili in $(x_{0},\xi_{k-1})$, e
$G^{(k)}\neq0$ in $(x_{0},\xi_{k-1})$. Per il teorema di Cauchy (\cref{teo:cauchy}) esiste
$\xi_{k}\in(x_{0},\xi_{k-1})$ con
\[
  \frac{F^{(k)}(\xi_{k})}{G^{(k)}(\xi_{k})}=\frac{F^{(k-1)}(\xi_{k-1})-F^{(k-1)}(x_{0})}{G^{(k-1)}(\xi_{k-1})-G^{(k-1)}(x_{0})}
  =\frac{F^{(k-1)}(\xi_{k-1})}{G^{(k-1)}(\xi_{k-1})} ,
\]
perché $F^{(k-1)}(x_{0})=G^{(k-1)}(x_{0})=0$. Concatenando le \eqref{eq:catena-cauchy} e posto
$x_{1}=\xi_{n+1}\in(x_{0},x)$,
\[
  \frac{R_{n}(x)}{(x-x_{0})^{n+1}}=\frac{F(\xi_{0})}{G(\xi_{0})}=\frac{F^{(n+1)}(x_{1})}{G^{(n+1)}(x_{1})}
  =\frac{f^{(n+1)}(x_{1})}{(n+1)!} ,
\]
che è la tesi.
\end{dimostrazione}

\begin{strategia}[Un'altra via]
Invece di derivare $n+1$ volte, teniamo fisso il punto $x$ e facciamo variare il
\emph{centro} dello sviluppo. La funzione $\Phi(t)$, valore in $x$ del polinomio di Taylor
con centro $t$, vale $T_{n}(x)$ per $t=x_{0}$ e $f(x)$ per $t=x$; la sua derivata, per una
cancellazione «telescopica», contiene soltanto $f^{(n+1)}$. Basta allora \emph{una sola}
applicazione del teorema di Cauchy, confrontando $\Phi$ con $\Psi(t)=(x-t)^{n+1}$: una funzione
ausiliaria ben scelta sostituisce la catena di passaggi della prima dimostrazione.
\end{strategia}

\begin{dimostrazione}[Seconda dimostrazione]
Per $t\in I$ poniamo
\begin{equation}\label{eq:taylor-centro-variabile}
  \Phi(t)=\sum_{k=0}^{n}\frac{f^{(k)}(t)}{k!}\,(x-t)^{k}, \qquad \Psi(t)=(x-t)^{n+1} .
\end{equation}
Per $t=x_{0}$ la somma è il polinomio di Taylor con centro $x_{0}$ calcolato in $x$, e per
$t=x$ sopravvive soltanto il termine con $k=0$; dunque $\Phi(x_{0})=T_{n}(x)$, $\Phi(x)=f(x)$ e
$\Phi(x)-\Phi(x_{0})=R_{n}(x)$, mentre $\Psi(x)-\Psi(x_{0})=-(x-x_{0})^{n+1}$. Poiché $f^{(k)}$ è
derivabile in $I$ per $k\le n$, la funzione $\Phi$ è derivabile in $I$, e per la regola del
prodotto
\[
\begin{aligned}
  \Phi'(t)&=\sum_{k=0}^{n}\frac{f^{(k+1)}(t)}{k!}(x-t)^{k}
  -\sum_{k=1}^{n}\frac{f^{(k)}(t)}{(k-1)!}(x-t)^{k-1}\\
  &=\frac{f^{(n+1)}(t)}{n!}\,(x-t)^{n} :
\end{aligned}
\]
posto $j=k-1$, la seconda somma diventa $\sum_{j=0}^{n-1}\frac{f^{(j+1)}(t)}{j!}(x-t)^{j}$, che si elide
con i termini della prima di indice $k\le n-1$, e resta soltanto quello con $k=n$. Inoltre
$\Psi'(t)=-(n+1)(x-t)^{n}\neq0$ per $t\neq x$. Le funzioni $\Phi$ e $\Psi$ sono continue e derivabili
nell'intervallo chiuso di estremi $x_{0}$ e $x$, contenuto in $I$, e $\Psi'$ non si annulla
nei suoi punti interni; poiché il rapporto degli incrementi non cambia scambiando gli
estremi, il teorema di Cauchy (\cref{teo:cauchy}) fornisce un punto $x_{1}$ strettamente
compreso tra $x_{0}$ e $x$ con
\[
  \frac{R_{n}(x)}{-(x-x_{0})^{n+1}}=\frac{\Phi'(x_{1})}{\Psi'(x_{1})}=-\frac{f^{(n+1)}(x_{1})}{(n+1)!} ,
\]
dove si è semplificato il fattore $(x-x_{1})^{n}\neq0$. È la \eqref{eq:resto-lagrange}.
\end{dimostrazione}

\begin{osservazione}[Le due dimostrazioni a confronto]\label{oss:confronto-resto-lagrange}
Entrambe le dimostrazioni usano soltanto il teorema di Cauchy, e per questo nessuna
delle due richiede la continuità di $f^{(n+1)}$. La prima applica il teorema $n+1$ volte,
a funzioni che si annullano in $x_{0}$ insieme alle loro derivate: è il teorema di
L'Hôpital «al finito». La seconda lo applica una sola volta, ma a una funzione costruita
con un'idea nuova, quella di trattare il centro dello sviluppo come variabile. La
funzione $\Phi$ ha anche un'altra virtù: la scelta di $\Psi$ è libera, e cambiando $\Psi$ si
ottengono altre forme del resto, come vedremo nel \S\ref{subsec:resto-cauchy}.
\end{osservazione}

In entrambe le versioni il resto di Lagrange paga, rispetto al resto integrale, un prezzo:
il punto $x_{1}$ è soltanto «garantito». Come nel teorema di Lagrange, non sappiamo dove si
trovi, se non che sta tra $x_{0}$ e $x$, e in generale dipende sia da $x$ sia da $n$. Per le stime, però, questo basta, come mostrano gli
esempi che seguono.

\begin{figure}[htb]
\centering
\begin{tikzpicture}
\begin{axis}[width=0.86\textwidth, height=6.2cm, xmin=-1.2, xmax=2.1, ymin=0, ymax=6.2,
  axis lines=middle, xlabel={$x$}, ylabel={$y$}, clip=false,
  xtick={0.4522,1.6}, xticklabels={$x_{1}$,$x$}, ytick=\empty]
  \addplot[curva, domain=-1.2:1.82, samples=80] {exp(x)};
  \addplot[curva2, tratteggio, domain=-1.2:2.05, samples=60] {1+x+x^2/2};
  \draw[verde, very thick] (axis cs:1.6,3.88) -- (axis cs:1.6,4.953);
  \fill[blu] (axis cs:1.6,4.953) circle (1.6pt);
  \fill[rame] (axis cs:1.6,3.88) circle (1.6pt);
  \draw[verde, thin] (axis cs:1.6,4.42) -- (axis cs:1.8,3.3);
  \node[verde, font=\scriptsize, anchor=north] at (axis cs:1.8,3.3) {$R_{2}(x)$};
  \fill (axis cs:0.4522,0) circle (1.4pt);
  \node[blu, font=\scriptsize, anchor=west] at (axis cs:1.84,6.05) {$e^{x}$};
  \node[rame, font=\scriptsize] at (axis cs:-0.9,1.7) {$T_{2}(x)$};
\end{axis}
\end{tikzpicture}
\caption{La funzione $e^{x}$ e il suo polinomio di Taylor $T_{2}(x)=1+x+\frac{x^{2}}2$ in $x_{0}=0$.
Per $x=1{,}6$ il resto vale $R_{2}(x)=e^{1{,}6}-3{,}88\approx1{,}073$; per il
\cref{teo:resto-lagrange} è $R_{2}(x)=\frac{e^{x_{1}}}{6}\,x^{3}$, e risolvendo si trova
$x_{1}=\log\frac{6R_{2}(x)}{x^{3}}\approx0{,}452$. Il punto $x_{1}$ non è il punto medio di
$[0,x]$.}
\label{fig:resto-lagrange}
\end{figure}

\begin{esempio}[Approssimazione dei piccoli angoli, con stima]\label{es:piccoli-angoli-stima}
Nel \cref{sec:taylor} abbiamo osservato che l'approssimazione $\sen\varphi\approx\varphi$ ha un errore
relativo di circa $\frac{\varphi^{2}}6$. Il resto di Lagrange trasforma l'«all'incirca» in una
disuguaglianza valida per ogni $\varphi$. Il polinomio di Taylor del seno di ordine $2$ in
$0$ è $T_{2}(\varphi)=\varphi$ (il coefficiente di $\varphi^{2}$ è nullo), e la derivata terza del seno è
$-\cos$; dunque, per qualche $\varphi_{1}$ tra $0$ e $\varphi$,
\[
  \sen\varphi-\varphi=-\frac{\cos\varphi_{1}}{6}\,\varphi^{3} ,
\]
e quindi
\begin{equation}\label{eq:stima-seno}
  \abs{\sen\varphi-\varphi}\le\frac{\abs\varphi^{3}}6 \qquad\perogni\varphi\in\R .
\end{equation}
Per $\varphi=10^{\circ}\approx0{,}1745$~rad l'errore relativo $\frac{\abs{\sen\varphi-\varphi}}{\abs{\sen\varphi}}$ è quindi al
più $\frac{\varphi^{3}/6}{\varphi-\varphi^{3}/6}\approx0{,}51\%$. Si noti l'utilità di aver scelto $n=2$ invece di
$n=1$: i due polinomi coincidono, ma con $n=1$ si otterrebbe soltanto
$\abs{\sen\varphi-\varphi}\le\frac{\varphi^{2}}2$, una stima molto peggiore per $\varphi$ piccolo.
\end{esempio}

\begin{python}[Il punto $x_{1}$ di Lagrange per $e^{x}$, con $x=1$]
import math
from scipy.integrate import quad

print(" n      R_n(1)       integrale       x_1     1/(n+2)")
for n in range(1, 8):
    T = sum(1 / math.factorial(k) for k in range(n + 1))
    R = math.e - T                                 # resto esatto
    I = quad(lambda t: (1 - t)**n / math.factorial(n)
             * math.exp(t), 0, 1)[0]               # resto integrale
    x1 = math.log(math.factorial(n + 1) * R)       # da e^x1/(n+1)!
    print(f"{n:2d}  {R:.9f}  {I:.9f}  {x1:.5f}  {1/(n+2):.5f}")
\end{python}
\begin{risultato}
 n      R_n(1)       integrale       x_1     1/(n+2)
 1  0.718281828  0.718281828  0.36225  0.33333
 2  0.218281828  0.218281828  0.26979  0.25000
 3  0.051615162  0.051615162  0.21411  0.20000
 4  0.009948495  0.009948495  0.17716  0.16667
 5  0.001615162  0.001615162  0.15093  0.14286
 6  0.000226273  0.000226273  0.13139  0.12500
 7  0.000027860  0.000027860  0.11629  0.11111
\end{risultato}

Il programma conferma che il resto integrale coincide con il resto calcolato per
differenza, e mostra che il punto $x_{1}$ si avvicina a $x_{0}=0$ al crescere di $n$, circa
come $\frac1{n+2}$. La ragione è che il peso $(1-t)^{n}$ si concentra sempre più vicino a $x_{0}$:
il suo baricentro in $[0,1]$ è
\[
  \frac{\int_{0}^{1}t\,(1-t)^{n}\,dt}{\int_{0}^{1}(1-t)^{n}\,dt}=\frac{\frac1{(n+1)(n+2)}}{\frac1{n+1}}=\frac1{n+2}
\]
(il numeratore si calcola con la sostituzione $u=1-t$), e la media pesata di $e^{t}$, che
vicino a $0$ è quasi lineare, viene raggiunta in un punto vicino al baricentro.

\subsection{Approfondimento: la forma di Cauchy del resto}
\label{subsec:resto-cauchy}

Il resto di Lagrange non è sempre la forma più conveniente. Sia $f(x)=\log(1+x)$, con
$x_{0}=0$ e $-1<x<0$. Poiché $f^{(n+1)}(t)=\frac{(-1)^{n}n!}{(1+t)^{n+1}}$ (\cref{sec:taylor}), la
\eqref{eq:resto-lagrange} dà, per qualche $x_{1}\in(x,0)$,
\[
  \abs{R_{n}(x)}=\frac{\abs x^{n+1}}{(n+1)(1+x_{1})^{n+1}}
  \le\frac1{n+1}\Bigl(\frac{\abs x}{1+x}\Bigr)^{n+1} ,
\]
perché $1+x_{1}>1+x$. Se $-\frac12\le x<0$ il rapporto $\frac{\abs x}{1+x}$ non supera $1$, e la stima
tende a zero per $n\to+\infty$; ma se $-1<x<-\frac12$ il rapporto supera $1$ e la stima
diverge. Poiché di $x_{1}$ sappiamo soltanto che sta in $(x,0)$, non possiamo fare di
meglio: il resto di Lagrange, da solo, non permette di decidere se $R_{n}(x)$ tenda a zero.
Una forma diversa del resto, dovuta a Cauchy, supera l'ostacolo.

\begin{teorema}[Forma di Cauchy del resto]\label{teo:resto-cauchy}
\index{formula di Taylor!con il resto di Cauchy}
Sia $f$ derivabile $n+1$ volte in un intervallo aperto $I$, e siano $x_{0},x\in I$ con
$x\neq x_{0}$. Esiste un punto $x_{2}$ strettamente compreso tra $x_{0}$ e $x$ tale che
\begin{equation}\label{eq:resto-cauchy}
  R_{n}(x)=\frac{f^{(n+1)}(x_{2})}{n!}\,(x-x_{2})^{n}\,(x-x_{0}) .
\end{equation}
\end{teorema}

\begin{schema}
\Ipotesi{$f$ derivabile $n+1$ volte in $I$ aperto; \ $x_{0}\neq x$ in $I$.}
\Tesi{$\esiste x_{2}$ strettamente compreso tra $x_{0}$ e $x$ per cui vale la
\eqref{eq:resto-cauchy}.}
\end{schema}

\begin{strategia}
Si riprende la funzione $\Phi$ della \eqref{eq:taylor-centro-variabile}, per la quale
$\Phi(x)-\Phi(x_{0})=R_{n}(x)$ e $\Phi'(t)=\frac{f^{(n+1)}(t)}{n!}(x-t)^{n}$. Al teorema di Cauchy con
$\Psi(t)=(x-t)^{n+1}$ si sostituisce il teorema di Lagrange, cioè il teorema di Cauchy con
$\Psi(t)=t$: così la derivata $\Phi'$ viene calcolata \emph{per intero} nel punto intermedio,
compreso il fattore $(x-t)^{n}$.
\end{strategia}

\begin{dimostrazione}
Come nella seconda dimostrazione del \cref{teo:resto-lagrange-generale}, la funzione $\Phi$
è continua e derivabile nell'intervallo chiuso di estremi $x_{0}$ e $x$. Per il teorema di
Lagrange (\cref{teo:lagrange}), applicato in questo intervallo (il rapporto
$\frac{\Phi(x)-\Phi(x_{0})}{x-x_{0}}$ non cambia scambiando gli estremi), esiste $x_{2}$ strettamente
compreso tra $x_{0}$ e $x$ con
\[
\begin{aligned}
  R_{n}(x)&=\Phi(x)-\Phi(x_{0})=\Phi'(x_{2})\,(x-x_{0})\\
  &=\frac{f^{(n+1)}(x_{2})}{n!}\,(x-x_{2})^{n}\,(x-x_{0}) .
\end{aligned}
\qedhere
\]
\end{dimostrazione}

Torniamo al logaritmo. Per $-1<x<x_{2}<0$ la \eqref{eq:resto-cauchy} dà
\[
  \abs{R_{n}(x)}=\frac{\abs x\,(x_{2}-x)^{n}}{(1+x_{2})^{n+1}}
  =\frac{\abs x}{1+x_{2}}\Bigl(\frac{x_{2}-x}{1+x_{2}}\Bigr)^{n}\le\frac{\abs x^{n+1}}{1+x} .
\]
Infatti $\frac1{1+x_{2}}<\frac1{1+x}$; e $\frac{x_{2}-x}{1+x_{2}}\le\abs x=-x$ equivale, moltiplicando per
$1+x_{2}>0$, a $x_{2}\le-x\,x_{2}$, cioè a $x_{2}(1+x)\le0$, che è vera. Ora la stima tende a zero
per ogni $x\in(-1,0)$. Il punto intermedio compare due volte nel fattore
$\frac{x_{2}-x}{1+x_{2}}$, e i due effetti si compensano: se $x_{2}$ è vicino a $x$, dove il
denominatore $1+x_{2}$ è piccolo, è piccolo anche il numeratore $x_{2}-x$.

\begin{osservazione}[Tre forme dello stesso resto]\label{oss:tre-forme-resto}
Se $f$ è di classe $C^{n+1}$, la forma di Cauchy si ricava anche dal resto integrale,
come fa Giusti \cite{giusti}: invece di portare fuori dall'integrale
\eqref{eq:resto-integrale} il solo fattore $f^{(n+1)}$, con il \cref{lem:media-pesata}, si
applica all'intero integrando $\frac{(x-t)^{n}}{n!}f^{(n+1)}(t)$ il teorema della media
(\cref{teo:media}). Le forme di Lagrange e di Cauchy sono dunque due «riassunti» diversi
della stessa formula integrale, e nessuna delle due è sempre la migliore. Più in
generale, con $\Psi(t)=\abs{x-t}^{p}$ e $p>0$ la seconda dimostrazione del
\cref{teo:resto-lagrange-generale} produce la forma del resto detta di Schlömilch e
Roche, che per $p=n+1$ e per $p=1$ restituisce le due precedenti. Quando si sa
stimare direttamente, il resto integrale resta il più preciso: per il logaritmo con
$-1<x<0$ il \cref{sec:taylor-peano} ha dato $\abs{R_{n}(x)}\le\frac{\abs x^{n+1}}{(n+1)(1+x)}$, migliore
di un fattore $n+1$. Il programma confronta le tre stime con il resto vero.
\end{osservazione}

\begin{python}[Resto di $\log(1+x)$ in $x=-0{,}8$: tre stime]
import math

x = -0.8                          # log(1+x) con x < -1/2
q = abs(x) / (1 + x)              # = 4
print("  n    |R_n(x)|   Lagrange     Cauchy   integrale")
for n in (1, 2, 5, 10, 20, 40):
    T = sum((-1)**(k + 1) * x**k / k for k in range(1, n + 1))
    R = abs(math.log(1 + x) - T)
    lag = q**(n + 1) / (n + 1)
    cau = abs(x)**(n + 1) / (1 + x)
    itg = cau / (n + 1)
    print(f"{n:3d}  {R:.3e}  {lag:.3e}  {cau:.3e}  {itg:.3e}")
\end{python}
\begin{risultato}
  n    |R_n(x)|   Lagrange     Cauchy   integrale
  1  8.094e-01  8.000e+00  3.200e+00  1.600e+00
  2  4.894e-01  2.133e+01  2.560e+00  8.533e-01
  5  1.508e-01  6.827e+02  1.311e+00  2.185e-01
 10  3.056e-02  3.813e+05  4.295e-01  3.905e-02
 20  1.892e-03  2.094e+11  4.612e-02  2.196e-03
 40  1.192e-05  1.179e+23  5.317e-04  1.297e-05
\end{risultato}

Per $x=-0{,}8$ il rapporto $\frac{\abs x}{1+x}$ vale $4$: la stima di Lagrange cresce come
$\frac{4^{n+1}}{n+1}$ e non dice nulla, mentre quella di Cauchy decresce come $\frac{0{,}8^{n+1}}{0{,}2}$ e
quella integrale segue da vicino il resto vero.

\begin{notastorica}[Lagrange, Cauchy e il resto]
Per oltre mezzo secolo il teorema di Taylor rimase un risultato tra tanti, finché
Lagrange, in una memoria del 1772 per l'Accademia di Berlino, non lo proclamò il
principio fondamentale del calcolo differenziale. Nella \emph{Théorie des fonctions
analytiques} (1797), il cui frontespizio è l'epigrafe di questo capitolo, tentò di
fondare l'intero calcolo sugli sviluppi in serie, definendo la derivata come
coefficiente del termine di primo grado, senza infinitesimi né limiti (si veda la nota
storica del \S\ref{sec:def-derivata}). Interessato ai calcoli approssimati, Lagrange volle anche sapere
quale errore si commette arrestando la serie, e trovò il resto del
\cref{teo:resto-lagrange}: lo strinse tra i valori minimo e massimo di $f^{(n+1)}$ e
concluse con la proprietà dei valori intermedi, che allora si riteneva evidente, secondo
lo stesso schema del \cref{lem:media-pesata}. Il programma conteneva però un circolo
vizioso, come osserva Giusti \cite{giusti}: per sapere se una funzione è uguale alla
propria serie bisogna studiarne il resto, e quindi proprio le derivate che si volevano
evitare. Cauchy, nel \emph{Résumé des leçons} (1823), rovesciò l'ordine: definita la
derivata come limite, la formula di Taylor diventò un teorema, e il resto, che egli
scrisse anche nella forma \eqref{eq:resto-cauchy}, ne divenne il contenuto essenziale.
Cauchy mostrò inoltre, con la funzione $e^{-1/x^{2}}$ (\cref{prop:cauchy-controesempio}), che
funzioni diverse possono avere la stessa serie di Taylor: una funzione non è determinata
dalla sua serie, e la definizione di Lagrange non poteva reggere. Del suo programma
sopravvissero il resto e la notazione $f'$.
\end{notastorica}
